\section{Systementwurf}
\subsection{User Interface}
Zur Planung der Benutzeroberfläche wurden im Vorfeld Wireframes und Mockups mit dem Design-Tool Figma erstellt. 
Diese dienen dazu, die Struktur der einzelnen Seiten frühzeitig zu visualisieren und zu überprüfen. 
Interaktive Funktionalitäten oder Navigationsflüsse wurden im Rahmen dieses Projekts nicht prototypisch umgesetzt. 
Die Entwürfe umfassen unter anderem:
\begin{itemize}
	\item die Karriereseite (vgl. \ref{fig:landingPage})
	\item die Stellenübersicht (vgl. \ref{fig:stellenOverview})
	\item die detailierte Stellenanzeige (vgl. \ref{fig:stelle})
	\item ein Bewerbungsformular für Kandidat:innen (vgl. \ref{fig:bewerbungFormular})
	\item ein Dashboard fuer Recruiter:innen (vgl. \ref{fig:recruiterDashboard})
\end{itemize}
Die weiteren erstellten WireFrames bzw. Mockups koennen unter folgenden 
\href{https://www.figma.com/design/0LDsbppkYlQgEbpulk5ud9/ATS-WebApp?node-id=20-46&t=cHVE2zw3fKD8hl14}{\textit{Link}} eingesehen werden.

\subsection{Datenmodell}
Datenmodell wurde mithilfe eines ER-Modell visualisiert.
Das ERM ist im Anhang unter --Hier Abbildung ERM referenzieren-- zu finden.

\subsection{Technologieauswahl}
Aufgrund der Vergleiche, die in den Tabellen \ref{tab:vglBackend} und \ref{tab:dbvergleich} durchgefuehert wurden, ergbibt sich folgender Tech-Stack.
Das Backend wird mittels PHP, das Frontend via dem CSS Framework Tailwind implementiert und als Datenbank wird das leichgewichtige SQLite DBMS genutzt.

\subsection{Infrastruktur}
Die Infrastruktur des Systems ist in drei Umgebungen unterteilt: Entwicklungs-, Test- und Produktivumgebung. 
Die Entwicklung erfolgt auf der lokalen Entwicklermaschine, die Testumgebung wird mittels VM gehostet und der prototypische Betrieb der Anwendung erfolgt auf einem Raspberry Pi. 
Auf dem Gerät läuft eine Linux-basierte Serverumgebung. :w
Der zentrale Apache-Web-Server und die PHP-Anwendung werden mithilfe des Tools Docker containisiert.
Durch die Nutzung von Docker sind die Umgebungen standardisiert, was einen Rollout von Test auf Prod relativ leicht gestaltet.
Der PI ist fuer diesen MVP vollkommend ausreichend. Falls das System spaeter hochskaliert werden soll, koennen Cloud- oder Virtual-Server eine Rolle spielen.
Ein vollstaendiges Architekturbild ist im Anhang unter 

\subsection{Softwarearchitektur}
Für die Umsetzung des Bewerbermanagementsystems wird eine klassische 3-Schichten-Architektur gewählt. 
Diese unterteilt die Anwendung in Präsentationsschicht, Anwendungslogik und Datenhaltung. 
Die Präsentationsschicht bildet die Benutzeroberfläche für Bewerber:innen und Recruiter:innen ab, während die Geschäftslogik in PHP umgesetzt wird und die Verarbeitung von Bewerbungen, Statuswechseln und Verwaltungsfunktionen steuert. 
Die Datenhaltung erfolgt über eine SQLite-Datenbank, in der Stellenanzeigen, Bewerbungen, Nutzer und Dokumente strukturiert gespeichert werden.
Die gewählte Architektur eignet sich besonders für den prototypischen Charakter des Systems, da sie eine klare Trennung der Verantwortlichkeiten ermöglicht und dadurch Wartbarkeit, Übersichtlichkeit und Erweiterbarkeit verbessert. 
Gleichzeitig bleibt die technische Umsetzung schlank und angemessen für den Umfang eines Minimum Viable Product.
